% =============================================================
%  Void Linux + Sway — دليل شامل
%  Compile with: xelatex void-sway-guide.tex
% =============================================================
\documentclass[12pt, a4paper]{report}

% ── Language & Font ──────────────────────────────────────────
\usepackage{fontspec}
\usepackage{polyglossia}
\setmainlanguage[numerals=mashriq]{arabic}
\setotherlanguage{english}

\setmainfont{Noto Serif Arabic}[Script=Arabic, Scale=1.1]
\newfontfamily\arabicfont{Noto Sans Arabic}[Script=Arabic]
\newfontfamily\englishfont{JetBrains Mono}[Scale=0.9]
\setmonofont{JetBrains Mono}[Scale=0.85]

% ── Page Layout ──────────────────────────────────────────────
\usepackage[
  a4paper,
  top=2.5cm, bottom=2.5cm,
  left=2.8cm, right=2.8cm
]{geometry}

% ── Colors ───────────────────────────────────────────────────
\usepackage[dvipsnames, table]{xcolor}
\definecolor{codebg}{HTML}{1E1E2E}
\definecolor{codefg}{HTML}{CDD6F4}
\definecolor{accent}{HTML}{89B4FA}
\definecolor{warn}{HTML}{FAB387}
\definecolor{success}{HTML}{A6E3A1}
\definecolor{tabhead}{HTML}{313244}
\definecolor{notecolor}{HTML}{F5C2E7}
\definecolor{rulecolor}{HTML}{6C7086}

% ── Code Listings ────────────────────────────────────────────
\usepackage{listings}
\lstdefinestyle{bashstyle}{
  language=bash,
  basicstyle=\englishfont\small\color{codefg},
  backgroundcolor=\color{codebg},
  frame=single,
  framesep=6pt,
  rulecolor=\color{rulecolor},
  breaklines=true,
  breakatwhitespace=true,
  showstringspaces=false,
  tabsize=2,
  xleftmargin=10pt,
  xrightmargin=10pt,
  keywordstyle=\color{accent}\bfseries,
  commentstyle=\color{success}\itshape,
  stringstyle=\color{warn},
  moredelim=[is][\color{notecolor}]{@}{@},
}
\lstdefinestyle{inistyle}{
  basicstyle=\englishfont\small\color{codefg},
  backgroundcolor=\color{codebg},
  frame=single,
  framesep=6pt,
  rulecolor=\color{rulecolor},
  breaklines=true,
  xleftmargin=10pt,
  xrightmargin=10pt,
}
\lstset{style=bashstyle}

% ── Tables ───────────────────────────────────────────────────
\usepackage{booktabs}
\usepackage{tabularx}
\usepackage{array}
\renewcommand{\arraystretch}{1.4}
\newcolumntype{R}{>{\raggedleft\arraybackslash}X}
\newcolumntype{L}{>{\raggedright\arraybackslash}X}
\newcolumntype{C}{>{\centering\arraybackslash}X}

% ── Misc Packages ────────────────────────────────────────────
\usepackage{graphicx}
\usepackage{hyperref}
\usepackage{enumitem}
\usepackage{titlesec}
\usepackage{fancyhdr}
\usepackage{tcolorbox}
\tcbuselibrary{skins, breakable}
\usepackage{mdframed}
\usepackage{microtype}
\usepackage{multicol}

% ── Hyperref Setup ───────────────────────────────────────────
\hypersetup{
  colorlinks=true,
  linkcolor=accent,
  urlcolor=accent,
  pdftitle={Void Linux + Sway — دليل شامل},
  pdfauthor={دليل شخصي},
}

% ── Custom Boxes ─────────────────────────────────────────────
\tcbset{
  warningbox/.style={
    enhanced, breakable,
    colback=codebg, colframe=warn,
    fonttitle=\bfseries\color{warn},
    title={⚠️\ تحذير},
    left=6pt, right=6pt, top=4pt, bottom=4pt,
    arc=3pt,
  },
  notebox/.style={
    enhanced, breakable,
    colback=codebg, colframe=accent,
    fonttitle=\bfseries\color{accent},
    title={ℹ\ ملاحظة},
    left=6pt, right=6pt, top=4pt, bottom=4pt,
    arc=3pt,
  },
}

% ── Section Formatting ───────────────────────────────────────
\titleformat{\chapter}[display]
  {\normalfont\huge\bfseries\color{accent}}
  {\chaptertitlename\ \thechapter}{10pt}{\Huge}
\titleformat{\section}
  {\normalfont\Large\bfseries\color{accent!80!white}}
  {\thesection}{1em}{}
\titleformat{\subsection}
  {\normalfont\large\bfseries\color{white!80!black}}
  {\thesubsection}{1em}{}

% ── Header / Footer ──────────────────────────────────────────
\pagestyle{fancy}
\fancyhf{}
\fancyhead[R]{\small\color{rulecolor}Void Linux + Sway}
\fancyhead[L]{\small\color{rulecolor}\leftmark}
\fancyfoot[C]{\small\color{rulecolor}\thepage}
\renewcommand{\headrulewidth}{0.3pt}

% ── Helper commands ──────────────────────────────────────────
\newcommand{\pkg}[1]{\texttt{\color{accent}#1}}
\newcommand{\cmd}[1]{\texttt{\color{codefg}#1}}
\newcommand{\file}[1]{\texttt{\color{warn}#1}}

% =============================================================
\begin{document}
\pagecolor{codebg}
\color{codefg}

% ─── Title Page ──────────────────────────────────────────────
\begin{titlepage}
  \centering
  \vspace*{3cm}
  {\Huge\bfseries\color{accent} Void Linux + Sway\par}
  \vspace{1cm}
  {\huge\color{white} دليل شامل\par}
  \vspace{1cm}
  \rule{0.5\textwidth}{0.5pt}
  \vspace{1cm}
  \par
  {\large\color{codefg!80!white}
  دليل شخصي لإعادة بناء البيئة من الصفر على Void Linux + Sway.\\[0.4em]
  مبني على تجربة فعلية والدوكس الرسمية.\\[0.4em]
  كل فصل مستقل --- افتح اللي تحتاجه بس.
  \par}
  \vspace{2cm}
  {\large\color{rulecolor} الترتيب المنطقي للتثبيت من الصفر:\par}
  \vspace{0.5cm}
  {\large\color{accent!80!white}
  01 → 02 → 03 → 04 → 05 → 06 → 07\par}
  \vfill
  \rule{0.4\textwidth}{0.3pt}
\end{titlepage}

% ─── TOC ─────────────────────────────────────────────────────
\tableofcontents
\clearpage

% ─── Quick Reference ─────────────────────────────────────────
\chapter*{قواعد Void الأساسية}
\addcontentsline{toc}{chapter}{قواعد Void الأساسية}

\begin{lstlisting}
# تثبيت برنامج
sudo xbps-install PACKAGE

# تفعيل service
sudo ln -s /etc/sv/SERVICENAME /var/service/

# بعد أي usermod -aG --> لازم logout وادخل تاني

# تشغيل سواي الصح
dbus-run-session sway

# elogind vs seatd --> اختار واحد بس، مش ممكن الاتنين
\end{lstlisting}

% =============================================================
\chapter{التثبيت الأساسي}
% =============================================================

\section{تحديث النظام أولاً}

\begin{lstlisting}
sudo xbps-install -Su
\end{lstlisting}

\section{أدوات البناء الأساسية}

\begin{lstlisting}
sudo xbps-install base-devel git make ca-certificates wget curl
\end{lstlisting}

\section{GPU Drivers (Mesa بدون Xorg)}

\begin{tcolorbox}[notebox]
على Wayland مش محتاج \texttt{xorg-server} أو \texttt{xf86-video-intel}.
\end{tcolorbox}

\subsection{Intel (مثل Dell Latitude E6430)}
\begin{lstlisting}
sudo xbps-install mesa-dri intel-video-accel \
  libva-utils mesa-vaapi mesa-vdpau \
  vulkan-loader mesa-vulkan-intel
\end{lstlisting}

\subsection{AMD}
\begin{lstlisting}
sudo xbps-install mesa-dri mesa-vaapi mesa-vdpau \
  vulkan-loader mesa-vulkan-radeon
\end{lstlisting}

\subsection{NVIDIA}
\begin{lstlisting}
sudo xbps-install nvidia nvidia-libs
\end{lstlisting}

\section{Session Manager (اختار واحد بس!)}

\begin{tcolorbox}[warningbox]
elogind و seatd مش بيشتغلوا مع بعض --- اختار واحد بس.
\end{tcolorbox}

\subsection{الخيار A: elogind (الأنسب للاب)}
\begin{lstlisting}
sudo xbps-install elogind
sudo ln -s /etc/sv/elogind /var/service/
\end{lstlisting}

\textbf{مميزات elogind:}
\begin{itemize}[leftmargin=*]
  \item بيعمل \texttt{XDG\_RUNTIME\_DIR} تلقائياً
  \item بيدعم power management (suspend، hibernate)
  \item أكثر استقراراً مع معظم البرامج
\end{itemize}

\subsection{الخيار B: seatd (أخف وأبسط)}
\begin{lstlisting}
sudo xbps-install seatd
sudo ln -s /etc/sv/seatd /var/service/
sudo usermod -aG _seatd $USER
\end{lstlisting}

\textbf{مميزات seatd:}
\begin{itemize}[leftmargin=*]
  \item أخف بكتير
  \item بيدير الـ seat بس (مش XDG\_RUNTIME\_DIR)
  \item محتاج \texttt{pam\_rundir} معاه
\end{itemize}

\section{XDG\_RUNTIME\_DIR}

\subsection{لو اخترت elogind}
مش محتاج حاجة --- بيتعمل تلقائياً.

\subsection{لو اخترت seatd}
\begin{lstlisting}
sudo xbps-install pam_rundir
\end{lstlisting}
أضف في \file{/etc/pam.d/system-login}:
\begin{lstlisting}[style=inistyle]
-session optional pam_rundir.so
\end{lstlisting}

\section{D-Bus}
\begin{lstlisting}
sudo xbps-install dbus
sudo ln -s /etc/sv/dbus /var/service/
\end{lstlisting}

\section{المجلدات الافتراضية}
\begin{lstlisting}
sudo xbps-install xdg-user-dirs
xdg-user-dirs-update
\end{lstlisting}

\section{Sway والأدوات الأساسية}
\begin{lstlisting}
sudo xbps-install \
  sway swaybg swayidle swaylock \
  xorg-server-xwayland \
  wlroots wl-clipboard

# Terminal
sudo xbps-install tilix   # او foot (اخف واسرع)
sudo xbps-install foot

# App Launcher
sudo xbps-install fuzzel   # او wofi
sudo xbps-install wofi

# Status Bar
sudo xbps-install waybar

# Notifications
sudo xbps-install mako libnotify

# Screenshot
sudo xbps-install grim slurp

# Screen Recording
sudo xbps-install wf-recorder

# Clipboard Manager
sudo xbps-install clipman
\end{lstlisting}

\section{XDG Desktop Portals (ضروري لـ Screen Sharing)}
\begin{lstlisting}
sudo xbps-install \
  xdg-desktop-portal \
  xdg-desktop-portal-wlr \
  xdg-desktop-portal-gtk \
  xdg-utils
\end{lstlisting}

\subsection{ملف إعداد الـ Portals (مهم جداً!)}
\begin{lstlisting}
mkdir -p ~/.config/xdg-desktop-portal
nano ~/.config/xdg-desktop-portal/sway-portals.conf
\end{lstlisting}

\begin{lstlisting}[style=inistyle]
[preferred]
default=gtk
org.freedesktop.impl.portal.Screenshot=wlr
org.freedesktop.impl.portal.ScreenCast=wlr
\end{lstlisting}

\begin{tcolorbox}[notebox]
اسم الملف لازم يكون \texttt{sway-portals.conf} بالظبط لأن \texttt{XDG\_CURRENT\_DESKTOP=sway}.
\end{tcolorbox}

\section{Media Codecs}
\begin{lstlisting}
sudo xbps-install \
  ffmpeg \
  gst-plugins-base1 \
  gst-plugins-good1 \
  gst-plugins-bad1 \
  gst-plugins-ugly1
\end{lstlisting}

\section{Power Management (للاب)}
\begin{lstlisting}
sudo xbps-install tlp acpid acpi brightnessctl
sudo ln -s /etc/sv/acpid /var/service/
sudo ln -s /etc/sv/tlpd /var/service/
sudo usermod -aG video $USER
\end{lstlisting}

\subsection{إعداد TLP (مهم لـ WiFi stability)}
\begin{lstlisting}
sudo nano /etc/tlp.conf
\end{lstlisting}
ابحث عن وغيّر:
\begin{lstlisting}[style=inistyle]
WIFI_PWR_ON_AC=off
WIFI_PWR_ON_BAT=off
\end{lstlisting}

\section{Fonts}
\begin{lstlisting}
sudo xbps-install \
  liberation-fonts-ttf \
  dejavu-fonts-ttf \
  noto-fonts-ttf \
  noto-fonts-ttf-extra \
  noto-fonts-emoji \
  font-ibm-plex-ttf \
  terminus-font \
  font-awesome6
\end{lstlisting}

\begin{tcolorbox}[notebox]
\textbf{JetBrains Mono:} مش في الريبو، تنزله يدوي من \href{https://www.jetbrains.com/lp/mono/}{jetbrains.com/lp/mono} وتحطه في \texttt{\textasciitilde/.local/share/fonts/}
\end{tcolorbox}

\section{متغيرات البيئة}

أضف في \file{\textasciitilde/.bash\_profile} أو \file{\textasciitilde/.profile}:

\begin{lstlisting}
# XDG Base Directories
export XDG_CONFIG_HOME="$HOME/.config"
export XDG_DATA_HOME="$HOME/.local/share"
export XDG_CACHE_HOME="$HOME/.cache"
export XDG_STATE_HOME="$HOME/.local/state"
export PATH="$HOME/.local/bin:$PATH"

# XDG_RUNTIME_DIR (لو مش شغّال elogind)
if test -z "${XDG_RUNTIME_DIR}"; then
    export XDG_RUNTIME_DIR=/tmp/${UID}-runtime-dir
    if ! test -d "${XDG_RUNTIME_DIR}"; then
        mkdir "${XDG_RUNTIME_DIR}"
        chmod 0700 "${XDG_RUNTIME_DIR}"
    fi
fi

# Wayland
export XDG_SESSION_TYPE=wayland
export XDG_SESSION_DESKTOP=sway
export XDG_CURRENT_DESKTOP=sway
export DE=sway

# Apps
export MOZ_ENABLE_WAYLAND=1
export QT_QPA_PLATFORM=wayland
export QT_WAYLAND_DISABLE_WINDOWDECORATION=1
export SDL_VIDEODRIVER=wayland
export ELECTRON_OZONE_PLATFORM_HINT=auto
export _JAVA_AWT_WM_NONREPARENTING=1

# لو بتستخدم seatd
export LIBSEAT_BACKEND=seatd

# Auto-start Sway من tty1
if [ -z "$DISPLAY" ] && [ -z "$WAYLAND_DISPLAY" ] && [ "${XDG_VTNR:-0}" -eq 1 ]; then
    exec dbus-run-session sway
fi
\end{lstlisting}

\section{نسخ الكونفيج الافتراضي}
\begin{lstlisting}
mkdir -p ~/.config/sway
cp /etc/sway/config ~/.config/sway/config
\end{lstlisting}

\section{إضافات أساسية في كونفيج سواي}

افتح \file{\textasciitilde/.config/sway/config} وأضف:

\begin{lstlisting}
### تحديث بيئة dbus ###
exec dbus-update-activation-environment DISPLAY WAYLAND_DISPLAY SWAYSOCK XDG_CURRENT_DESKTOP

### XDG Desktop Portal ###
exec /usr/libexec/xdg-desktop-portal-wlr &
exec /usr/libexec/xdg-desktop-portal -r &

### PipeWire ###
exec pipewire &

### Notifications ###
exec mako &

### Polkit ###
exec /usr/libexec/polkit-gnome-authentication-agent-1 &

### Network ###
exec nm-applet --indicator &

### Bluetooth ###
exec blueman-applet &

### Clipboard ###
exec wl-paste -t text --watch clipman store &

### USB Auto-mount ###
exec udiskie --tray &

### Keyboard Layout ###
input type:keyboard {
    xkb_layout us,ara
    xkb_options grp:lwin_toggle
    repeat_delay 300
    repeat_rate 50
}

### Touchpad ###
input type:touchpad {
    tap enabled
    natural_scroll enabled
    dwt enabled
}

### Shortcuts ###
# Screenshot
bindsym Print exec grim ~/Pictures/screenshot-$(date +%Y%m%d-%H%M%S).png
bindsym Shift+Print exec grim -g "$(slurp)" ~/Pictures/screenshot-$(date +%Y%m%d-%H%M%S).png
bindsym Ctrl+Print exec grim -g "$(slurp)" - | wl-copy

# Volume
bindsym XF86AudioRaiseVolume exec pactl set-sink-volume @DEFAULT_SINK@ +5%
bindsym XF86AudioLowerVolume exec pactl set-sink-volume @DEFAULT_SINK@ -5%
bindsym XF86AudioMute exec pactl set-sink-mute @DEFAULT_SINK@ toggle

# Brightness
bindsym XF86MonBrightnessUp exec brightnessctl set +10%
bindsym XF86MonBrightnessDown exec brightnessctl set 10%-

# Lock screen
bindsym $mod+l exec swaylock -f -c 000000

### Idle + Lock ###
exec swayidle -w \
    timeout 300 'swaylock -f -c 000000' \
    timeout 600 'swaymsg "output * dpms off"' \
    resume 'swaymsg "output * dpms on"' \
    before-sleep 'swaylock -f -c 000000'
\end{lstlisting}

\section{تحقق نهائي}
\begin{lstlisting}
sv status dbus
sv status elogind   # او seatd
groups $USER
echo $XDG_RUNTIME_DIR
echo $WAYLAND_DISPLAY   # بعد تشغيل سواي: wayland-1
pactl info | grep "Server Name"
grim ~/test.png
\end{lstlisting}

% =============================================================
\chapter{الصوت (PipeWire + ALSA)}
% =============================================================

\section{التثبيت}
\begin{lstlisting}
sudo xbps-install \
  pipewire \
  wireplumber \
  alsa-utils \
  alsa-tools \
  alsa-pipewire \
  libspa-bluetooth \
  pavucontrol
\end{lstlisting}

\begin{tcolorbox}[notebox]
\pkg{libspa-bluetooth} ضروري عشان الـ Bluetooth headphones تشتغل مع PipeWire.
\end{tcolorbox}

\section{إعداد WirePlumber}
\begin{lstlisting}
mkdir -p ~/.config/pipewire/pipewire.conf.d

ln -s /usr/share/examples/wireplumber/10-wireplumber.conf \
      ~/.config/pipewire/pipewire.conf.d/
\end{lstlisting}

\section{إعداد PipeWire-Pulse (لدعم برامج PulseAudio)}
\begin{lstlisting}
sudo xbps-install pipewire-pulse

ln -s /usr/share/examples/pipewire/20-pipewire-pulse.conf \
      ~/.config/pipewire/pipewire.conf.d/
\end{lstlisting}

\section{إعداد ALSA مع PipeWire}
\begin{lstlisting}
sudo mkdir -p /etc/alsa/conf.d

sudo ln -s /usr/share/alsa/alsa.conf.d/50-pipewire.conf \
           /etc/alsa/conf.d/

sudo ln -s /usr/share/alsa/alsa.conf.d/99-pipewire-default.conf \
           /etc/alsa/conf.d/

alsactl init
\end{lstlisting}

\section{التشغيل}

على Void، PipeWire \textbf{بيشتغل من داخل سواي}، مش كـ system service.

في كونفيج سواي:
\begin{lstlisting}
exec pipewire &
\end{lstlisting}

\begin{tcolorbox}[warningbox]
لا تحاول تفعّله بـ \texttt{ln -s /etc/sv/pipewire /var/service/} --- مش هيشتغل صح.
\end{tcolorbox}

\section{ضبط مستويات الصوت}
\begin{lstlisting}
# CLI
alsamixer

# GUI
pavucontrol
\end{lstlisting}

\section{التحقق}
\begin{lstlisting}
pactl info | grep "Server Name"
# المفروض: PulseAudio (on PipeWire ...)

pactl list sources short
pactl list sinks short
\end{lstlisting}

\section{مشاكل شائعة}

\textbf{الصوت مش شغّال:}
\begin{lstlisting}
pactl info
pipewire &
pactl info
\end{lstlisting}

\textbf{Bluetooth Headphones مش بتشتغل:}
\begin{lstlisting}
sudo xbps-install libspa-bluetooth
# بعدين restart سواي
\end{lstlisting}

\textbf{No Sound Device Found:}
\begin{lstlisting}
aplay -l
# لو فاضي، راجع GPU drivers والـ kernel modules
\end{lstlisting}

% =============================================================
\chapter{الـ Theming (شكل موحد)}
% =============================================================

\section{المفهوم الأساسي}

الـ theme على Wayland = 4 طبقات منفصلة:

\begin{center}
\begin{tabularx}{0.85\textwidth}{>{\bfseries\color{accent}}l X}
\toprule
\rowcolor{tabhead} \color{white}الطبقة & \color{white}التأثير \\
\midrule
GTK Theme  & ألوان وشكل الـ buttons والـ windows \\
Icon Theme & أيقونات الملفات والبرامج \\
Cursor Theme & شكل المؤشر \\
Font & الخط في كل البرامج \\
\bottomrule
\end{tabularx}
\end{center}

\begin{tcolorbox}[notebox]
\textbf{المشكلة:} على Wayland، GTK مش بيقرأ الإعدادات تلقائياً --- لازم تتحط في \texttt{gsettings}.\\
\textbf{الحل:} أداة \pkg{nwg-look}.
\end{tcolorbox}

\section{التثبيت}
\begin{lstlisting}
sudo xbps-install nwg-look

# Materia (flat, modern)
sudo xbps-install materia-theme

# Arc
sudo xbps-install arc-theme

# Adwaita:dark -- موجود تلقائياً مع GTK
\end{lstlisting}

ثيم من النت (مثال Catppuccin):
\begin{lstlisting}
mkdir -p ~/.local/share/themes
cd ~/.local/share/themes
# نزّل وفكّ الـ zip هنا
\end{lstlisting}

\begin{lstlisting}
# Icon Theme
sudo xbps-install papirus-icon-theme

# Cursor Theme
sudo xbps-install breeze-cursors
\end{lstlisting}

\section{التطبيق}
\begin{lstlisting}
nwg-look
# اختار GTK Theme + Icons + Cursor + Font
# اضغط Apply
\end{lstlisting}

\section{Cursor في كونفيج سواي}

أضف في \file{\textasciitilde/.config/sway/config}:
\begin{lstlisting}
seat seat0 xcursor_theme اسم_الثيم 24
\end{lstlisting}

\section{الفونتات}
\begin{lstlisting}
sudo xbps-install \
  font-jetbrains-mono \
  noto-fonts-ttf \
  noto-fonts-ttf-extra \
  noto-fonts-emoji \
  font-awesome6 \
  dejavu-fonts-ttf

# تحديث cache بعد أي font جديد
fc-cache -fv

# دعم العربية
sudo xbps-install font-arabic-misc noto-fonts-ttf
\end{lstlisting}

في كونفيج سواي:
\begin{lstlisting}
font pango:JetBrains Mono 11, Noto Sans Arabic 11, Font Awesome 6 Free 11
\end{lstlisting}

\section{مكان الثيمات}
\begin{lstlisting}
# GTK themes
/usr/share/themes/
~/.local/share/themes/

# Icons + Cursors
/usr/share/icons/
~/.local/share/icons/

# شوف المثبت
ls /usr/share/themes/
ls /usr/share/icons/
fc-list | grep -i "jet\|noto\|arabic"
\end{lstlisting}

% =============================================================
\chapter{Screenshot + Screen Recording + Clipboard}
% =============================================================

\section{الأدوات}
\begin{lstlisting}
sudo xbps-install grim slurp wf-recorder wl-clipboard clipman
\end{lstlisting}

\section{Screenshot}

\subsection{أوامر مباشرة}
\begin{lstlisting}
# شاشة كاملة
grim ~/Pictures/screenshot-$(date +%Y%m%d-%H%M%S).png

# منطقة معينة
grim -g "$(slurp)" ~/Pictures/screenshot-$(date +%Y%m%d-%H%M%S).png

# نسخ مباشرة للـ clipboard
grim -g "$(slurp)" - | wl-copy
\end{lstlisting}

\subsection{في كونفيج سواي}
\begin{lstlisting}
bindsym Print exec grim ~/Pictures/screenshot-$(date +%Y%m%d-%H%M%S).png
bindsym Shift+Print exec grim -g "$(slurp)" ~/Pictures/screenshot-$(date +%Y%m%d-%H%M%S).png
bindsym Ctrl+Print exec grim -g "$(slurp)" - | wl-copy
\end{lstlisting}

\section{Screen Recording — wf-recorder}
\begin{lstlisting}
# شاشة كاملة
wf-recorder -f ~/Videos/recording.mp4

# منطقة معينة
wf-recorder -g "$(slurp)" -f ~/Videos/recording.mp4

# مع الصوت
wf-recorder -a -f ~/Videos/recording.mp4

# مع صوت device معين
pactl list sources short          # شوف الـ devices
wf-recorder -a "اسم_السورس" -f ~/Videos/recording.mp4

# إيقاف
# Ctrl+C
\end{lstlisting}

\section{Clipboard Manager — clipman}

\subsection{تشغيل (في كونفيج سواي)}
\begin{lstlisting}
exec wl-paste -t text --watch clipman store &
\end{lstlisting}

\subsection{عرض الـ history}
\begin{lstlisting}
clipman pick --tool wofi
# أو
clipman pick --tool fuzzel
\end{lstlisting}

أضف في كونفيج سواي:
\begin{lstlisting}
bindsym $mod+shift+v exec clipman pick --tool fuzzel
\end{lstlisting}

\subsection{مسح الـ history}
\begin{lstlisting}
clipman clear --all
\end{lstlisting}

% =============================================================
\chapter{الأجهزة (Bluetooth + Printer + USB + MTP)}
% =============================================================

\section{Bluetooth}

\subsection{التثبيت}
\begin{lstlisting}
sudo xbps-install bluez blueman libspa-bluetooth
sudo ln -s /etc/sv/bluetoothd /var/service/
sudo usermod -aG bluetooth $USER
# logout وادخل تاني
\end{lstlisting}

في كونفيج سواي:
\begin{lstlisting}
exec blueman-applet &
\end{lstlisting}

\subsection{الاستخدام (CLI)}
\begin{lstlisting}
bluetoothctl
power on
scan on
pair XX:XX:XX:XX:XX:XX
connect XX:XX:XX:XX:XX:XX
trust XX:XX:XX:XX:XX:XX
\end{lstlisting}

\subsection{التحقق}
\begin{lstlisting}
sv status bluetoothd
\end{lstlisting}

\section{Printer — CUPS}

\subsection{التثبيت}
\begin{lstlisting}
sudo xbps-install cups cups-filters gutenprint system-config-printer
# HP:
sudo xbps-install hplip
# Epson:
sudo xbps-install epson-inkjet-printer-escpr

sudo ln -s /etc/sv/cupsd /var/service/
sudo usermod -aG lpadmin $USER
\end{lstlisting}

\subsection{إضافة الطابعة}
افتح في المتصفح: \texttt{http://localhost:631}\\
ثم: Administration → Add Printer

\section{USB Automount — udiskie}

\subsection{التثبيت}
\begin{lstlisting}
sudo xbps-install udisks2 udiskie gvfs ntfs-3g exfat-utils
\end{lstlisting}

\subsection{Polkit Rule (عشان ما يطلبش password)}
\begin{lstlisting}
sudo mkdir -p /etc/polkit-1/rules.d/
sudo nano /etc/polkit-1/rules.d/50-udisks.rules
\end{lstlisting}

\begin{lstlisting}[language={}]
polkit.addRule(function(action, subject) {
  if (subject.isInGroup("wheel")) {
    if (action.id.startsWith("org.freedesktop.udisks2.")) {
      return polkit.Result.YES;
    }
  }
});
\end{lstlisting}

في كونفيج سواي:
\begin{lstlisting}
exec udiskie --tray &
\end{lstlisting}

\subsection{Unmount يدوي}
\begin{lstlisting}
udiskie-umount /run/media/$USER/اسم_الـdrive
udiskie-umount --all
\end{lstlisting}

\section{MTP (موبايل عبر USB)}
\begin{lstlisting}
sudo xbps-install go-mtpfs libmtp
sudo usermod -aG plugdev $USER
\end{lstlisting}

\subsection{إعداد FUSE}
\begin{lstlisting}
# شيل الـ # من السطر user_allow_other
sudo sed -i '/user_allow_other/s/^#//g' /etc/fuse.conf
\end{lstlisting}

\subsection{الاستخدام}
\begin{lstlisting}
# شوف الأجهزة المتصلة
mtp-detect

# Mount الجهاز
mkdir -p ~/phone
go-mtpfs ~/phone &

# Unmount
fusermount -u ~/phone
\end{lstlisting}

% =============================================================
\chapter{الشبكة (VPN + Hotspot + Tethering)}
% =============================================================

\section{إعداد NetworkManager}
\begin{lstlisting}
sudo xbps-install NetworkManager network-manager-applet

# إيقاف الخدمات القديمة
sudo sv stop dhcpcd wpa_supplicant 2>/dev/null
sudo rm /var/service/dhcpcd /var/service/wpa_supplicant 2>/dev/null

# تفعيل NetworkManager
sudo ln -s /etc/sv/dbus /var/service/
sudo ln -s /etc/sv/NetworkManager /var/service/
\end{lstlisting}

\subsection{DNS ثابت (اختياري)}
\begin{lstlisting}
sudo nano /etc/resolv.conf
\end{lstlisting}
\begin{lstlisting}[style=inistyle]
nameserver 1.1.1.1
nameserver 8.8.8.8
\end{lstlisting}
\begin{lstlisting}
sudo chattr +i /etc/resolv.conf   # يمنع NetworkManager من تغييره
\end{lstlisting}

\section{VPN}

\subsection{تثبيت الـ plugins}
\begin{lstlisting}
# OpenVPN
sudo xbps-install NetworkManager-openvpn

# WireGuard
sudo xbps-install wireguard-tools NetworkManager-wireguard

# L2TP/IPSec
sudo xbps-install NetworkManager-l2tp
\end{lstlisting}

\subsection{إضافة VPN}
\begin{lstlisting}
# Import ملف config
nmcli connection import type openvpn file /path/to/config.ovpn
nmcli connection import type wireguard file /path/to/config.conf

# بـ TUI
nmtui
\end{lstlisting}

\subsection{تشغيل/إيقاف}
\begin{lstlisting}
nmcli connection up اسم_الـVPN
nmcli connection down اسم_الـVPN
nmcli connection show
\end{lstlisting}

\section{Hotspot}

\begin{tcolorbox}[warningbox]
معظم كروت WiFi مش بتدعم إنترنت + hotspot على نفس الكارت.
\end{tcolorbox}

\subsection{تشغيل سريع}
\begin{lstlisting}
nmcli device wifi hotspot ifname wlan0 ssid "اسم_الشبكة" password "كلمة_السر"
\end{lstlisting}

\subsection{عرض الـ QR Code}
\begin{lstlisting}
nmcli device wifi show-password
\end{lstlisting}

\subsection{Hotspot دائم}
\begin{lstlisting}
nmcli con add type wifi ifname wlan0 con-name "MyHotspot" autoconnect no \
  ssid "اسم_شبكتك" \
  802-11-wireless.mode ap \
  802-11-wireless.band bg \
  ipv4.method shared \
  wifi-sec.key-mgmt wpa-psk \
  wifi-sec.psk "كلمة_السر"

nmcli con up MyHotspot
nmcli con down MyHotspot
\end{lstlisting}

\section{USB Tethering (إنترنت الموبايل على اللاب)}
\begin{lstlisting}
# وصّل الموبايل بـ USB وفعّل USB Tethering منه
ip link show             # هيظهر interface جديد (usb0 او enp...)
nmcli device status      # NetworkManager المفروض يشوفه تلقائي
\end{lstlisting}

% =============================================================
\chapter{البرامج}
% =============================================================

\section{متصفح}
\begin{lstlisting}
sudo xbps-install firefox firefox-i18n-ar
\end{lstlisting}
Firefox على Wayland بيشتغل تلقائياً لو \texttt{MOZ\_ENABLE\_WAYLAND=1} موجود في البيئة.

\section{Multimedia}
\begin{lstlisting}
# Video
sudo xbps-install mpv        # native Wayland
# أو
sudo xbps-install vlc

# Images
sudo xbps-install imv         # native Wayland

# Audio control
sudo xbps-install pavucontrol
\end{lstlisting}

\section{File Manager}
\begin{lstlisting}
# Terminal
sudo xbps-install nnn

# GUI
sudo xbps-install thunar
\end{lstlisting}

\section{Archive}
\begin{lstlisting}
sudo xbps-install unzip zip p7zip file-roller
\end{lstlisting}

\section{System Monitoring}
\begin{lstlisting}
sudo xbps-install htop btop fastfetch lm_sensors
\end{lstlisting}

\section{Flatpak}
\begin{lstlisting}
sudo xbps-install flatpak xdg-user-dirs
flatpak remote-add --if-not-exists flathub https://dl.flathub.org/repo/flathub.flatpakrepo

# تثبيت
flatpak install flathub com.discordapp.Discord

# تشغيل
flatpak run com.discordapp.Discord

# تحديث الكل
flatpak update

# قائمة المثبت
flatpak list
\end{lstlisting}

\section{AppImage}
\begin{lstlisting}
sudo xbps-install fuse
sudo modprobe fuse

chmod +x اسم_البرنامج.AppImage
./اسم_البرنامج.AppImage

# لو بيشتكي من FUSE
./اسم_البرنامج.AppImage --appimage-extract-and-run

# مكان منظم
mkdir -p ~/AppImages
\end{lstlisting}

\section{Emacs}
\begin{lstlisting}
sudo xbps-install emacs         # مع GUI
sudo xbps-install emacs-nox     # terminal فقط (أخف)
\end{lstlisting}
الكونفيج في: \file{\textasciitilde/.config/emacs/init.el}

\section{Terminal Tools للـ Coding}
\begin{lstlisting}
sudo xbps-install \
  ripgrep \
  fd \
  bat \
  eza \
  fzf \
  zoxide \
  btop \
  lazygit
\end{lstlisting}

\begin{center}
\begin{tabularx}{0.9\textwidth}{>{\ttfamily\color{accent}}l >{\ttfamily}l X}
\toprule
\rowcolor{tabhead}\color{white}الأداة & \color{white}بديل عن & \color{white}الميزة \\
\midrule
rg      & grep  & أسرع بكتير \\
fd      & find  & أسهل وأسرع \\
bat     & cat   & syntax highlighting \\
eza     & ls    & أجمل وأوضح \\
fzf     & ---   & fuzzy finder لكل حاجة \\
zoxide  & cd    & بيتذكر الأماكن \\
btop    & htop  & أجمل وأكتر معلومات \\
lazygit & git   & TUI لـ git \\
\bottomrule
\end{tabularx}
\end{center}

\subsection{ربط fzf مع الـ history}
أضف في \file{\textasciitilde/.bashrc}:
\begin{lstlisting}
source /usr/share/fzf/key-bindings.bash
source /usr/share/fzf/completion.bash
\end{lstlisting}
بعدين \texttt{Ctrl+R} هيفتح fuzzy search في الـ history.

% =============================================================
\chapter{AppImage + Desktop Entries}
% =============================================================

\section{هيكل المجلدات}
\begin{lstlisting}[style=inistyle]
~/AppImages/
 +-- obsidian.AppImage
 +-- ticktick.AppImage
 +-- icons/
     +-- obsidian.png
     +-- ticktick.png

~/.local/share/applications/    # Desktop entries
~/.config/mimeapps.list         # MIME associations
\end{lstlisting}

\section{إعداد AppImage جديد}
\begin{lstlisting}
# 1. أنشئ المجلدات
mkdir -p ~/AppImages/icons
mkdir -p ~/.local/share/applications

# 2. انقل الـ AppImage
mv downloaded.AppImage ~/AppImages/appname.AppImage
chmod 755 ~/AppImages/appname.AppImage

# 3. حط الأيقونة
cp icon.png ~/AppImages/icons/appname.png
chmod 644 ~/AppImages/icons/appname.png
\end{lstlisting}

\section{Desktop Entry Template}
\begin{lstlisting}
nano ~/.local/share/applications/appname.desktop
\end{lstlisting}

\begin{lstlisting}[style=inistyle]
[Desktop Entry]
Version=1.0
Type=Application
Name=اسم_البرنامج
Comment=وصف قصير
Exec=/home/USERNAME/AppImages/appname.AppImage %F
Icon=/home/USERNAME/AppImages/icons/appname.png
Terminal=false
Categories=Office;Utility;
MimeType=text/markdown;text/plain;
StartupNotify=true
\end{lstlisting}

\begin{lstlisting}
chmod 644 ~/.local/share/applications/appname.desktop
update-desktop-database ~/.local/share/applications
\end{lstlisting}

\section{Exec Variables}

\begin{center}
\begin{tabularx}{0.7\textwidth}{>{\ttfamily\color{accent}}l X}
\toprule
\rowcolor{tabhead}\color{white}المتغير & \color{white}المعنى \\
\midrule
\%f & ملف واحد \\
\%F & عدة ملفات \\
\%u & URL واحد \\
\%U & عدة URLs \\
\bottomrule
\end{tabularx}
\end{center}

\section{أهم MIME Types}
\begin{lstlisting}[style=inistyle]
text/plain        ->  .txt
text/markdown     ->  .md
text/html         ->  .html
application/pdf   ->  .pdf
image/png         ->  .png
image/jpeg        ->  .jpg
\end{lstlisting}

\subsection{أوامر مفيدة}
\begin{lstlisting}
# اعرف MIME type لملف
xdg-mime query filetype document.md

# اعرف التطبيق الافتراضي
xdg-mime query default text/markdown

# حدّد التطبيق الافتراضي
xdg-mime default obsidian.desktop text/markdown
\end{lstlisting}

\section{مثال كامل: Obsidian}
\begin{lstlisting}
mkdir -p ~/AppImages/icons
mv Obsidian-*.AppImage ~/AppImages/obsidian.AppImage
chmod 755 ~/AppImages/obsidian.AppImage
cp obsidian-icon.png ~/AppImages/icons/obsidian.png
chmod 644 ~/AppImages/icons/obsidian.png
\end{lstlisting}

\begin{lstlisting}[style=inistyle]
# ~/.local/share/applications/obsidian.desktop
[Desktop Entry]
Version=1.0
Type=Application
Name=Obsidian
Comment=Knowledge base and note-taking app
Exec=/home/USERNAME/AppImages/obsidian.AppImage %F
Icon=/home/USERNAME/AppImages/icons/obsidian.png
Terminal=false
Categories=Office;TextEditor;Utility;
MimeType=text/markdown;text/plain;
StartupNotify=true
\end{lstlisting}

\begin{lstlisting}
chmod 644 ~/.local/share/applications/obsidian.desktop
update-desktop-database ~/.local/share/applications
xdg-mime default obsidian.desktop text/markdown
\end{lstlisting}

\section{Web App كـ Desktop Entry}
\begin{lstlisting}[style=inistyle]
[Desktop Entry]
Version=1.0
Type=Application
Name=Gmail
Exec=firefox --new-window https://mail.google.com
Icon=/home/USERNAME/AppImages/icons/gmail.png
Terminal=false
Categories=Network;WebBrowser;
StartupNotify=true
\end{lstlisting}

\section{Quick Reference}
\begin{lstlisting}
chmod +x file.AppImage
chmod 644 file.desktop
update-desktop-database ~/.local/share/applications
xdg-mime default app.desktop text/markdown
xdg-mime query default text/markdown
gtk-launch appname          # تجربة الـ desktop entry
xdg-open file.md            # فتح بالتطبيق الافتراضي
desktop-file-validate app.desktop
\end{lstlisting}

% =============================================================
\chapter{KVM/QEMU Virtualization}
% =============================================================

\section{التأكد من دعم CPU للـ Virtualization}
\begin{lstlisting}
lscpu | grep -i virtualization
grep -E '(vmx|svm)' /proc/cpuinfo
\end{lstlisting}

\begin{itemize}[leftmargin=*]
  \item \textbf{Intel:} لازم يظهر \texttt{VT-x} أو \texttt{vmx}
  \item \textbf{AMD:} لازم يظهر \texttt{AMD-V} أو \texttt{svm}
\end{itemize}

لو مش ظاهر، فعّله من الـ BIOS: ابحث عن \textit{Virtualization Technology} أو \textit{Intel VT-x}.

\section{التثبيت}
\begin{lstlisting}
sudo xbps-install virt-manager libvirt qemu dnsmasq bridge-utils openbsd-netcat
\end{lstlisting}

\begin{center}
\begin{tabularx}{0.85\textwidth}{>{\ttfamily\color{accent}}l X}
\toprule
\rowcolor{tabhead}\color{white}الحزمة & \color{white}الوظيفة \\
\midrule
virt-manager  & GUI لإدارة الـ VMs \\
libvirt       & الـ virtualization API \\
qemu          & المحاكي \\
dnsmasq       & DNS/DHCP للشبكات الافتراضية \\
bridge-utils  & network bridging \\
\bottomrule
\end{tabularx}
\end{center}

\section{تفعيل الـ Services}
\begin{lstlisting}
sudo ln -s /etc/sv/dbus /var/service/        # لو مش متفعّل
sudo ln -s /etc/sv/libvirtd /var/service/
sudo ln -s /etc/sv/virtlockd /var/service/
sudo ln -s /etc/sv/virtlogd /var/service/
\end{lstlisting}

\section{إضافة المستخدم لـ Group}
\begin{lstlisting}
sudo usermod -aG libvirt $USER
# logout وادخل تاني
\end{lstlisting}

\section{تحميل KVM Module}
\begin{lstlisting}
# Intel
sudo modprobe kvm_intel

# AMD
sudo modprobe kvm_amd

# تحميل تلقائي عند الـ Boot (Intel)
echo "kvm_intel" | sudo tee /etc/modules-load.d/kvm.conf
\end{lstlisting}

\section{إعداد libvirt (لتشغيل VMs بدون sudo)}
\begin{lstlisting}
sudo nano /etc/libvirt/libvirtd.conf
\end{lstlisting}

شيل الـ \# من:
\begin{lstlisting}[style=inistyle]
unix_sock_group = "libvirt"
unix_sock_ro_perms = "0777"
unix_sock_rw_perms = "0770"
\end{lstlisting}

\begin{lstlisting}
sudo nano /etc/libvirt/libvirt.conf
\end{lstlisting}

شيل الـ \# من:
\begin{lstlisting}[style=inistyle]
uri_default = "qemu:///system"
\end{lstlisting}

\begin{lstlisting}
sudo sv restart libvirtd
sudo reboot
\end{lstlisting}

\section{التحقق بعد الـ Reboot}
\begin{lstlisting}
lsmod | grep kvm
ls -la /dev/kvm
sudo virsh list --all
groups | grep libvirt
\end{lstlisting}

\section{أوامر virsh المهمة}
\begin{lstlisting}
virsh list --all              # كل الـ VMs
virsh start <vm-name>         # تشغيل
virsh shutdown <vm-name>      # إيقاف ناعم
virsh destroy <vm-name>       # إيقاف قسري
virsh undefine <vm-name>      # حذف
virsh dominfo <vm-name>       # معلومات
virsh edit <vm-name>          # تعديل الكونفيج
\end{lstlisting}

\section{Nested Virtualization}
\begin{lstlisting}
# Intel
sudo modprobe -r kvm_intel
sudo modprobe kvm_intel nested=1

# دائم
echo "options kvm_intel nested=1" | sudo tee /etc/modprobe.d/kvm.conf
\end{lstlisting}

\section{Troubleshooting}
\begin{lstlisting}
# "Unable to connect to libvirt"
sudo sv status dbus
sudo sv status libvirtd
sudo sv up libvirtd

# "Could not access KVM kernel module"
lsmod | grep kvm
sudo modprobe kvm_intel
ls -la /dev/kvm

# "Permission denied"
groups | grep libvirt
sudo usermod -aG libvirt $USER
# logout وادخل تاني
\end{lstlisting}

% =============================================================
\chapter{Troubleshooting}
% =============================================================

\section{سواي مش بيشتغل}

\subsection{"Failed to create backend"}
\begin{lstlisting}
# تأكد من elogind أو seatd
sudo sv status elogind
sudo sv status seatd

# لو seatd، تأكد من الـ group
groups | grep _seatd

# تأكد من dbus
sudo sv status dbus
\end{lstlisting}

\subsection{سواي بيبدأ وبيوقف بسرعة}
\begin{lstlisting}
dbus-run-session sway 2>&1 | head -30
\end{lstlisting}

\section{XDG\_RUNTIME\_DIR}
\begin{lstlisting}
# حل يدوي مؤقت
export XDG_RUNTIME_DIR=/run/user/$(id -u)
sudo mkdir -p $XDG_RUNTIME_DIR
sudo chown $USER:$USER $XDG_RUNTIME_DIR
sudo chmod 0700 $XDG_RUNTIME_DIR

# حل دائم: pam_rundir (مع seatd)
sudo xbps-install pam_rundir
# أضف في /etc/pam.d/system-login:
# -session optional pam_rundir.so
\end{lstlisting}

\section{الصوت}

\subsection{PipeWire مش شغّال}
\begin{lstlisting}
sudo sv status dbus
ls ~/.config/pipewire/pipewire.conf.d/
pipewire &
pactl info
grep pipewire ~/.config/sway/config
\end{lstlisting}

\section{Firefox مش شغّال على Wayland}
\begin{lstlisting}
echo $MOZ_ENABLE_WAYLAND
# المفروض: 1

MOZ_ENABLE_WAYLAND=1 firefox
\end{lstlisting}

\section{Screen Sharing مش شغّالة}
\begin{lstlisting}
# 1. تأكد من التثبيت
xbps-query xdg-desktop-portal-wlr

# 2. تأكد من portals config
cat ~/.config/xdg-desktop-portal/sway-portals.conf

# 3. تأكد من متغيرات الـ dbus
dbus-update-activation-environment --print-environment | grep XDG_CURRENT_DESKTOP

# 4. Restart portals يدوياً
killall xdg-desktop-portal xdg-desktop-portal-wlr 2>/dev/null
/usr/libexec/xdg-desktop-portal-wlr &
/usr/libexec/xdg-desktop-portal -r &

# 5. اختبار في Firefox
# https://mozilla.github.io/webrtc-landing/gum_test.html
\end{lstlisting}

\section{Arabic text مش بيظهر صح}
\begin{lstlisting}
sudo xbps-install font-arabic-misc noto-fonts-ttf
fc-cache -fv
# في Tilix: Preferences -> Profile -> Custom Font -> Noto Sans Arabic
\end{lstlisting}

\section{برامج Qt بتظهر X11}
\begin{lstlisting}
echo $QT_QPA_PLATFORM
# المفروض: wayland

sudo xbps-install qt5-wayland qt6-wayland
\end{lstlisting}

\section{Java Apps (شاشة رمادية فارغة)}
\begin{lstlisting}
# أضف في ~/.profile
export _JAVA_AWT_WM_NONREPARENTING=1
\end{lstlisting}

\section{AppImage مش بيشتغل}
\begin{lstlisting}
sudo xbps-install fuse
sudo modprobe fuse
chmod 755 ~/AppImages/app.AppImage
./app.AppImage --appimage-extract-and-run
\end{lstlisting}

\section{أوامر تشخيص عامة}
\begin{lstlisting}
echo $WAYLAND_DISPLAY        # المفروض: wayland-1
swaymsg -t get_version
swaymsg -t get_outputs
swaymsg -t get_inputs
dbus-update-activation-environment --print-environment
\end{lstlisting}

% =============================================================
\chapter{xbps-src: بناء البرامج اللي مش في الريبو}
% =============================================================

\section{الفكرة الأساسية}

\pkg{xbps-src} = زي \texttt{PKGBUILD} على Arch أو \texttt{ebuild} على Gentoo.

بدل ما تبني من سورس يدوياً، بتكتب (أو بتنسخ) "وصفة" صغيرة اسمها \textbf{template}، وxbps-src يبني البرنامج ويعمله \texttt{.xbps} package تقدر تثبته وتشيله بـ xbps عادي.

\begin{tcolorbox}[notebox]
\textbf{الفايدة الكبيرة:} البرنامج يبقى managed --- تقدر تشيله بـ \texttt{xbps-remove} زي أي package تاني.
\end{tcolorbox}

\section{الإعداد الأولي (مرة واحدة بس)}

\begin{lstlisting}
sudo xbps-install git xtools base-devel

cd ~
git clone --depth=1 https://github.com/void-linux/void-packages.git
cd void-packages

./xbps-src binary-bootstrap
\end{lstlisting}

\begin{tcolorbox}[notebox]
\texttt{-{}-depth=1} بتجيب آخر snapshot بس من غير كل الـ history --- أسرع بكتير.
\end{tcolorbox}

\section{بناء Package موجود في void-packages}

\begin{lstlisting}
# ابحث عن الـ template
ls srcpkgs/ | grep -i اسم_البرنامج
./xbps-src search اسم_البرنامج

# بناء
./xbps-src pkg اسم_البرنامج

# التثبيت
sudo xbps-install --repository=hostdir/binpkgs اسم_البرنامج

# لو في nonfree
sudo xbps-install --repository=hostdir/binpkgs/nonfree اسم_البرنامج
\end{lstlisting}

\subsection{مثال عملي: Google Chrome}
\begin{lstlisting}
cd ~/void-packages
echo "XBPS_ALLOW_RESTRICTED=yes" >> etc/conf
./xbps-src pkg google-chrome
sudo xbps-install --repository=hostdir/binpkgs/nonfree google-chrome
\end{lstlisting}

\section{إضافة الـ local repo بشكل دائم}
\begin{lstlisting}
sudo nano /etc/xbps.d/local-repo.conf
\end{lstlisting}

\begin{lstlisting}[style=inistyle]
repository=/home/اسمك/void-packages/hostdir/binpkgs
repository=/home/اسمك/void-packages/hostdir/binpkgs/nonfree
\end{lstlisting}

\section{تشريح الـ Template}

\begin{lstlisting}
pkgname=my-tool
version=1.2.3
revision=1
short_desc="أداة رائعة"
maintainer="اسمك <email@example.com>"
license="MIT"
homepage="https://github.com/someone/my-tool"

distfiles="https://github.com/someone/my-tool/archive/v${version}.tar.gz"
checksum="sha256 hash هنا"

hostmakedepends="python3"
depends="python3"

do_install() {
    python3 setup.py install --prefix=/usr --root="${DESTDIR}"
}
\end{lstlisting}

\subsection{جدول أهم المتغيرات}

\begin{center}
\begin{tabularx}{0.9\textwidth}{>{\ttfamily\color{accent}}l X}
\toprule
\rowcolor{tabhead}\color{white}المتغير & \color{white}المعنى \\
\midrule
pkgname        & اسم الـ package \\
version        & الإصدار \\
revision       & رقم revision (ابدأ بـ 1) \\
short\_desc    & وصف قصير (مش أكثر من 72 حرف) \\
distfiles      & رابط السورس (tar.gz أو zip) \\
checksum       & الـ sha256 hash للملف \\
hostmakedepends & أدوات البناء \\
makedepends    & مكتبات وقت البناء \\
depends        & مكتبات وقت التشغيل \\
build\_style   & طريقة البناء التلقائية \\
\bottomrule
\end{tabularx}
\end{center}

\subsection{الـ build\_style}
\begin{lstlisting}
build_style=cmake          # CMake projects
build_style=meson          # Meson projects
build_style=gnu-configure  # Autotools (./configure)
build_style=python3-module # Python (setup.py)
build_style=go             # Go
build_style=cargo          # Rust
build_style=gnu-makefile   # برامج بسيطة (make فقط)
\end{lstlisting}

\section{الحصول على الـ checksum}
\begin{lstlisting}
# xtools
xgensum -i srcpkgs/اسم_البرنامج/template

# يدوي
wget https://رابط_السورس
sha256sum اسم_الملف.tar.gz
\end{lstlisting}

\section{مثال كامل: برنامج Rust (zoxide)}
\begin{lstlisting}
./xbps-src newpkg zoxide
\end{lstlisting}

\begin{lstlisting}[style=inistyle]
# srcpkgs/zoxide/template
pkgname=zoxide
version=0.9.4
revision=1
build_style=cargo
short_desc="Smarter cd command"
maintainer="انا <me@example.com>"
license="MIT"
homepage="https://github.com/ajeetdsouza/zoxide"
distfiles="https://github.com/ajeetdsouza/zoxide/archive/v${version}.tar.gz"
checksum="ضع sha256 هنا"
\end{lstlisting}

\begin{lstlisting}
xgensum -i srcpkgs/zoxide/template
./xbps-src pkg zoxide
sudo xbps-install --repository=hostdir/binpkgs zoxide
\end{lstlisting}

\section{أوامر xbps-src اليومية}
\begin{lstlisting}
./xbps-src pkg اسم_البرنامج          # بناء
./xbps-src -f pkg اسم_البرنامج       # بناء من الأول
./xbps-src clean اسم_البرنامج        # تنظيف
./xbps-src show اسم_البرنامج         # شوف dependencies
cd ~/void-packages && git pull        # تحديث void-packages
./xbps-src update-sys                 # تحديث packages قديمة
\end{lstlisting}

\section{تحديث package}
\begin{lstlisting}
cd ~/void-packages
git pull
nano srcpkgs/اسم_البرنامج/template   # غيّر version=
xgensum -i srcpkgs/اسم_البرنامج/template
./xbps-src pkg اسم_البرنامج
sudo xbps-install -u اسم_البرنامج
\end{lstlisting}

\section{Hold: منع تحديث package}
\begin{lstlisting}
sudo xbps-pkgdb -m hold اسم_البرنامج
sudo xbps-pkgdb -m unhold اسم_البرنامج
\end{lstlisting}

\section{Troubleshooting xbps-src}

\textbf{"checksum mismatch":}
\begin{lstlisting}
xgensum -i srcpkgs/اسم_البرنامج/template
\end{lstlisting}

\textbf{بناء بطيء جداً:}
\begin{lstlisting}
echo "XBPS_MAKEJOBS=$(nproc)" >> etc/conf
\end{lstlisting}

\textbf{شوف الـ log كامل:}
\begin{lstlisting}
./xbps-src pkg اسم_البرنامج 2>&1 | tee /tmp/build.log
less /tmp/build.log
\end{lstlisting}

\section{ملخص سريع}
\begin{lstlisting}
# إعداد أولي (مرة واحدة)
git clone --depth=1 https://github.com/void-linux/void-packages.git
cd void-packages
./xbps-src binary-bootstrap

# تثبيت package موجود
./xbps-src pkg اسم_البرنامج
sudo xbps-install --repository=hostdir/binpkgs اسم_البرنامج

# package جديد
./xbps-src newpkg اسم_البرنامج
# عدّل template
xgensum -i srcpkgs/اسم_البرنامج/template
./xbps-src pkg اسم_البرنامج
sudo xbps-install --repository=hostdir/binpkgs اسم_البرنامج

# إزالة
sudo xbps-remove اسم_البرنامج
\end{lstlisting}

% =============================================================
\end{document}
